\subsection{The primary trace and the single clearing integer}
\label{subsec:ftb-primary-trace-clearing}

We retain the integral model, its canonical formal frame and the
same-model comparisons specified above.  In this subsection \(V\)
is a complete discrete valuation ring of characteristic zero with
finite residue field of characteristic \(p\), \(F=\operatorname{Frac}V\),
and \(v\) is its additive valuation.  All inequalities may equally
be divided by any fixed positive normalization integer.
The chart condition \(\mathrm{Good3}\) means that its formal
truncations identify the formal group law and that its PEL
uniformization identifies torsion, including its polarized pairing.
The connection bound used below is the all-line-power inequality
\(\mathrm{ChBound}(\rho)\): every integral harmonic chart input of
the prescribed row degrees has value of valuation at least
\(-\rho\).  For the principal host it is
\eqref{eq:ftb-host-chart-loss}.

Put \(R=\mathcal O_E\), \(R^n_{\rm st}=R^n\subset E^n\).
A lattice \(\Lambda\subset E^n\) is \emph{unimodular} here precisely
when
\[
 \Lambda=\{x:\ \sum_j\bar x_jy_j\in R\text{ for every }y\in\Lambda\}.
\]
Write \(\Lambda_0=\Lambda\cap R^n_{\rm st}\).
The level conditions are that \(N\) is odd, \(M\mid N\), and
\(NR^n_{\rm st}\subset\Lambda_0\).
The condition \(\operatorname{TrivAt}(\Lambda,m)\) means that there
exists \(k\) prime to \(m\) with
\(k\Lambda\subset R^n_{\rm st}\) and
\(kR^n_{\rm st}\subset\Lambda\).
For \(r=p_{\rm sig}+q_{\rm sig}\), put
\[
 I_N=[\Lambda_0:NR^n_{\rm st}]^r,\qquad
 J_N=[\Lambda:NR^n_{\rm st}]^r .
\]
These are the source's \srcid{mIdx} and \srcid{zCard}.
The letter \(p\) in the residue characteristic is unrelated to
\(p_{\rm sig}\).

For a lattice \(L\), let \(\det L\) denote its determinant fractional
ideal, generated by determinants of \(n\) vectors in \(L\).
For the coefficient field \(K\), write
\[
  T(L)=\prod_{\tau:E\hookrightarrow\mathbb C}
               (\tau(\det L)\mathcal O_K)^{r_\tau}.
\]
Thus the small scalar of Lemma~\ref{lem:ftb-reflex-small-scalar}
satisfies
\begin{equation}\label{eq:ftb-small-scalar-ideal}
                         (s)T(\Lambda_0)\subset T(\Lambda).
\end{equation}
The normalization scalar \(\beta_\Lambda\) satisfies
\[
              (\beta_\Lambda)T(\Lambda)=T(\Lambda_0).
\]
Only this equality and its nonvanishing are needed in the local
argument; its complex absolute values were used in the norm bound.

\begin{srlemma}[Adapted matrices at a prime]
\label{lem:ftb-adapted-matrices}
Assume that \(E/E^+\) is unramified at every finite prime.
For a unimodular \(\Lambda\) with the level conditions above, there
are integers \(k,k_1\) prime to \(2p\), a unitary
\(\gamma\in\operatorname{GL}_n(E)\), and matrices
\(P,Q,G\in M_n(R)\), such that
\[
 \begin{gathered}
 k\gamma\Lambda\subset R^n_{\rm st},\qquad
 kR^n_{\rm st}\subset\gamma\Lambda,\qquad
 GR^n_{\rm st}\subset\Lambda_0,\qquad
 k_1\Lambda_0\subset GR^n_{\rm st},\\
 GQ=k_1N1,\qquad
                         PQ=kN\gamma .
 \end{gathered}
\]
In particular \((PQ)^*(PQ)=(kN)^2 1\).
For every valuation of \(K\) above \(p\), with the evident embeddings,
\begin{align}
 v_p\!\left(|N_{E/\mathbb Q}\det Q|^r\right)&=v_p(I_N),
                                                \label{eq:ftb-index-cancel}\\
 v(T(\Lambda_0))&=v\!\left(\prod_\tau\tau(\det G)^{r_\tau}\right),\\
 v(T(\Lambda))&=-v\!\left(\prod_\tau\tau(\det\gamma)^{r_\tau}\right).
\end{align}
Consequently \eqref{eq:ftb-small-scalar-ideal} implies
\begin{equation}\label{eq:ftb-adapted-scalar-integral}
                  v\!\left(s\prod_\tau\tau(\det P)^{r_\tau}\right)
                         \ge0 .
\end{equation}
\end{srlemma}
\begin{proof}
We spell out the local unitary step, including the dyadic case.
Over each completion of \(E^+\), its integral quadratic algebra is
either split or unramified quadratic.
A unimodular Hermitian lattice over it has an orthonormal basis.
In the split case this follows by choosing dual bases for the two
factors.
In the unramified field case, reduction is a nondegenerate Hermitian
space over a quadratic finite-field extension.
It contains a vector of nonzero norm: otherwise polarization of
\(h(x+ay,x+ay)\), as \(a\) varies, would force
\(\operatorname{Tr}(a h(x,y))=0\) for every \(a\), contradicting
nondegeneracy.
Lift this vector; its norm is a unit.
The norm map on units of an unramified quadratic local extension is
surjective.  On residue units this is the finite-field norm; on
successive principal-unit quotients the correction is the
surjective residue trace.  Completeness lifts these corrections.
This proof works at residue characteristic two.
Scale the lifted vector to norm one, split its orthogonal
complement, and induct.

Choose these local unitary frames at the finitely many primes over
\(2p\).  The source obtains one rational unitary matrix by the
Cayley approximation argument.
First multiply each local frame by a norm-one scalar so that
\(1+U\) is invertible; the excluded scalars are the finitely many
roots of the determinant polynomial.
There are infinitely many integral norm-one scalars, for example
\((1-t\theta)/(1+t\theta)\), with
\(\bar\theta=-\theta\) and sufficiently small integral \(t\).
Then \(S=(1-U)(1+U)^{-1}\) is skew-Hermitian.
Approximate the finitely many local \(S\)'s sufficiently closely by
a rational skew-Hermitian matrix: approximate entries in \(E\) and
replace an approximant \(Y\) by the appropriately scaled
\(Y-Y^*\).
The inverse Cayley matrix
\(2(1+S)^{-1}-1\) is unitary.
Its denominator is nonzero because, at a complex embedding, a
skew-Hermitian matrix has no eigenvalue \(-1\).
A sufficiently close approximant carries the two completed
lattices onto one another.  More explicitly its relative matrix
has integral entries and is unitary, so its inverse, the adjoint,
also has integral entries.
Clearing the remaining denominators gives one integer \(k_2\)
prime to \(2p\) and a unitary \(\gamma\) with
\(k_2\gamma\Lambda\subset R^n\) and
\(k_2R^n\subset\gamma\Lambda\).
There is no bound on \(k_2\), and none is used.

Localize \(R\) by the rational integers prime to \(2p\).
It is a semilocal Dedekind domain and hence a principal ideal
domain: its nonzero prime ideals are precisely the finitely many
primes over \(2p\), and approximation supplies generators with
prescribed valuations there.
The localized \(\Lambda_0\) therefore has a basis.
Clear its denominators by integers prime to \(2p\), and obtain
integral columns \(G\) and \(k_1\) prime to \(2p\) with
\[
       GR^n\subset\Lambda_0,\qquad k_1\Lambda_0\subset GR^n .
\]
Since \(Ne_j\in\Lambda_0\), there is an integral \(Q\) with
\(GQ=k_1N1\).
The columns \(k_2\gamma G\) are integral, so set
\(P=k_2\gamma G\), \(k=k_1k_2\).
These are the required identities.

For the index calculation, consider
\[
 NR^n\subset\Lambda_0,\qquad
 GQR^n=k_1NR^n\subset GR^n\subset\Lambda_0 .
\]
The quotients
\(NR^n/GQR^n\) and \(\Lambda_0/GR^n\) are killed by \(k_1\);
their orders are prime to \(p\).
Multiplicativity of indices consequently shows that
\([\Lambda_0:NR^n]\) and
\([GR^n:GQR^n]=[R^n:QR^n]\) have equal \(p\)-adic valuation.
The last index is \(|N_{E/\mathbb Q}\det Q|\).
One may see this by restriction of scalars to the free
\(\mathbb Z\)-module \(R^n\): the determinant is the field norm of
\(\det Q\).  Raise to \(r\) to get \eqref{eq:ftb-index-cancel}.

At a valuation above \(p\), the same column arguments say that
\(\det G\) generates the local determinant ideal of \(\Lambda_0\)
and \((\det\gamma)^{-1}\) generates that of \(\Lambda\);
all factors \(k,k_1\) are units.
Moreover
\(\det P\cdot k_1^n=k^n\det\gamma\det G\).
Apply the embeddings with multiplicities \(r_\tau\) to this equality.
The inclusion \eqref{eq:ftb-small-scalar-ideal}, interpreted as an
inequality of local ideal valuations, is exactly
\eqref{eq:ftb-adapted-scalar-integral}.
For fractional ideals, the minimum valuation is attained by one
of their finitely many generators; products and powers add these
minimum valuations, which justifies this argument without choosing
global principal generators.
\end{proof}

We use additive notation for finite torsion groups in the next two
arguments.  For a matrix \(U\), its action on a tuple is the sum of
the corresponding endomorphism actions.
Let \(Z_\Lambda(N)\) be the group of \(N^2\)-torsion tuples \(z\)
annihilated by every integral matrix \(N\mu\) for which
\(\mu\Lambda\subset R^n\).
This is the source's \srcid{InZh} group, of cardinal \(J_N\).
On its \(p\)-primary part the adapted identities give
\begin{equation}\label{eq:ftb-primary-adapted-kernel}
                z\in Z_\Lambda(N)\quad\Longleftrightarrow\quad PQz=0 .
\end{equation}
Indeed \(k\gamma\) is one of its permitted lattice maps, so the
forward implication follows after cancelling the prime-to-\(p\)
factor \(k\).
Conversely, for every permitted \(\mu\), the inclusion
\(kR^n\subset\gamma\Lambda\) makes
\(k\mu\gamma^{-1}\) integral.
It gives an integral left factor of \(PQ\) equal to \(k^2N\mu\).
Cancelling \(k^2\) proves the required annihilation.
The adjoint equality for \(PQ\) also gives \(N^2z=0\), after
cancelling \(k^2\).
The same proof applies at level \(kN\).

\begin{srlemma}[From the formal trace to all primary torsion]
\label{lem:ftb-primary-trace}
Let \(\mu=PQ\) be an odd-level similitude with
\(\mu^*\mu=L^2 1\).
Assume that there are positive integers \(k_0,M_0,k_1\) and
\(G_0\in M_n(\mathcal O_E)\) such that
\[
 L^2=k_0M_0,\qquad (k_0,p)=(k_1,p)=1,\qquad
 G_0Q=k_1L1,
\]
and that all \(M_0\)-torsion is rational on the same generic fibre
(the source's \(\mathrm{AllRat}(M_0)\)).
Assume also that its \(p\)-primary kernel \(W\) is enumerated and
the product Weil pairing is trivial on \(W\times W\).
Let \(F=(F_b)\) be a reached tuple of sections of line power \(L^2\),
fixed by the \(p\)-primary \(L\)-torsion.
Let \(G_u\) be its normalized symmetric translate by \(u\in W\).
If \(t\,P(\mu X)\) is an integral harmonic polynomial, then
\[
 t\sum_{u\in W}\operatorname{hv}(G_u,P)
   =\#W_0\prod_\tau\tau(\det P)^{r_\tau}\,z,\qquad
                         v(z)\ge-\rho ,
 \quad W_0=W\cap\ker Q .
\]
The hypothesis that \(tP(\mu X)\) is integral includes its
prescribed harmonic row degrees.
\end{srlemma}
\begin{proof}
Let \(W_1\subset W\) be the subgroup whose points are small formal
points of the same integral chart.
It is a subgroup by the integral formal group law; inverses are
also integral, or may be written as positive powers in this finite
torsion group.
For each \(x\in W\), the reached section \(G_x\) is regular by the
reached-section result.
Apply the formal trace calculation to its normalized translations
by \(W_1\).
Finite height, split kernel and the normalized-unit
\(\kappa\)-ratio identify precisely this formal subgroup.
It gives
\[
 t\sum_{\alpha\in W_1}\operatorname{hv}(G_{\alpha+x},P)
       =\#(W_1\cap W_0)
             \prod_\tau\tau(\det P)^{r_\tau}\,z_x,
 \qquad v(z_x)\ge-\rho .
\]
The simultaneous translation may differ from the composition of the
chosen individual normalized translations.
Its \(b\)-th scalar ratio \(r_b\) satisfies \(r_b^{L^2}=1\), and
the square of \(\prod_b r_b\) is the product Weil commutator.
Isotropy gives \((\prod_b r_b)^2=1\).
Oddness then gives \(\prod_b r_b=1\), exactly.
Multilinearity of the harmonic value therefore removes these
ratios.  Each ratio is an integral unit, since a root of unity
is integral together with its inverse.

A point of \(W_0\) is killed by \(k_1L\), because the supplementary
matrix identity has that scalar product.
It is \(p\)-primary and \(k_1\) is prime to \(p\), so it is killed
by \(L\), and fixes \(F\) by hypothesis.
The preceding normalized-translation calculation implies
\(\operatorname{hv}(G_{x+w},P)=\operatorname{hv}(G_x,P)\)
for \(w\in W_0\).

Here is the exact finite-coset cancellation.
For a function \(f\) on a finite abelian \(W\) constant on
\(W_0\)-cosets, pass to \(W/W_0\).
The map \(W_1\to W/W_0\) has kernel \(W_1\cap W_0\).
Its image sums therefore satisfy
\[
 \sum_{\alpha\in W_1} f(\alpha+x)
   =\#(W_1\cap W_0)
       \sum_{h\in\operatorname{im}(W_1\to W/W_0)}\bar f(h+\bar x).
\]
Cancel this nonzero integer in \(F\), and sum over the quotient of
\(W/W_0\) by that image.
Finally the sum on \(W\) is \(\#W_0\) times its quotient sum.
Thus the required \(z\) is the sum of the corresponding \(z_x\).
The nonarchimedean inequality keeps \(v(z)\ge-\rho\).
This cancellation takes place as an equality in \(F\); it makes
no unsupported assertion that division by a group order preserves
integrality.
\end{proof}

\begin{srlemma}[Prime-to-residue translates, with the scalar]
\label{lem:ftb-coprime-translates-scalar}
Suppose the level-division tuple \(t_b\) has order prime to \(p\),
all \(N^2\)-torsion is rational, and \(P\in\operatorname{harmL}(\Lambda,a)\).
Let \(F_{z,t}\) be the two normalized translates at line power \(N^2\)
appearing in the theta value.  If
\eqref{eq:ftb-small-scalar-ideal} holds, then
\[
 v\!\left(\frac{s}{I_N N^{\deg P}}
           \sum_{z\in Z_\Lambda(N)}
                     \operatorname{hv}(F_{z,t},P)\right)\ge-\rho .
\]
\end{srlemma}
\begin{proof}
Decompose the finite group \(Z_\Lambda(N)=Z_p\times Z^{(p)}\)
into its primary and prime-to-primary factors.
This is a genuine product, obtained by Bezout projection on the
finite exponent; its primary factor has order
\(p^{v_p(J_N)}\).
For each fixed prime-to-\(p\) factor, compose its translation with
that by \(t\).
Normalized translations of coprime orders commute, since their
commutator has both orders and hence equals one.
The resulting base tuple is reached, and is fixed by the
\(p\)-primary \(N\)-torsion: multiplication pullback is fixed by
that kernel, and the coprime translation preserves this property.

Choose the adapted data of Lemma~\ref{lem:ftb-adapted-matrices}.
The lattice is unimodular and
\((kN)\mathcal O_E^n\subset\Lambda\), since
\(N\mathcal O_E^n\subset\Lambda\).
The exact lattice-pairing theorem
\srcid{zh\_isotropic\_and\_half\_isotropic}, applied at the odd
level \(kN\), therefore makes the product Weil pairing trivial
on these tuples; the adapted-kernel criterion puts them in
\(Z_\Lambda(kN)\).  Thus the isotropy hypothesis in the primary
trace lemma is supplied on this very group, with its actual
polarized pairing.
Pull the tuples back by the normalized multiplication-by-\(k\)
map, changing line power to \((kN)^2\).
It identifies their primary tuples by
\eqref{eq:ftb-primary-adapted-kernel}; multiplication by \(k\) is
an automorphism on primary torsion.
For the polynomial,
\[
                 (kN)^{-\deg P}P(PQX)=P(\gamma X).
\]
This has integral coefficients.
Indeed \(\operatorname{symO}(\Lambda)\) is the ring generated by
algebraic integral constants and
\(\sum_j\tau(b_j)X_{\tau,i,j}\), \(b\in\Lambda\).
Substitution by \(\gamma\) sends such a generator to the same
linear form with \(b\) replaced by \(\gamma b\).
Its coefficients are integral in \(V\) because \(k\gamma b\) is
integral and \(k\) is a unit.
Ring operations preserve this property.
Unitarity preserves the harmonic Laplacians, and column
substitution preserves row degrees and row invariance.
This is the full coefficient argument in
\srcid{MemRangeOfIsIntegral.map\_pushL\_mem\_FFV}.

For that application its level \(L=kN\) satisfies
\(L^2=k^2N^2\), with \((k^2,p)=1\) and all \(N^2\)-torsion
rational by hypothesis; its supplementary matrix is \(kG\),
since \((kG)Q=k_1(kN)1\).  Every supplementary hypothesis of
Lemma~\ref{lem:ftb-primary-trace} is now present.
It gives a scalar \(z\), \(v(z)\ge-\rho\).
Its kernel factor is
\[
  \#(Z_p\cap\ker Q)
      =p^{\,v_p(|N_{E/\mathbb Q}\det Q|^r)}
      =p^{v_p(I_N)} .
\]
For the first equality, identify the rational \(N^2\)-torsion with
the quotient of the integral PEL period lattice by \(N^2\).
Apply the source's primary-kernel count on \(N^2\)-torsion
with supplementary matrix \(NG\), since
\((NG)Q=k_1N^2 1\).
The resulting determinant-index calculation identifies the order
of the primary kernel with that of the primary part of
\(\operatorname{coker}Q\).
Both injectivity and surjectivity follow by cancelling \(k_1\)
using Bezout.
On the integral PEL lattice, which is projective of \(R\)-rank \(r\)
and free over \(\mathbb Z\), restriction of scalars gives
\[
            |\det_{\mathbb Z}Q|
                    =|N_{E/\mathbb Q}\det Q|^r .
\]
To avoid choosing an \(R\)-basis globally, choose a free
full-rank \(R\)-sublattice.  Its inclusion intertwines \(Q\);
after tensoring with \(\mathbb Q\) the two matrices are conjugate,
so have the same determinant.
The determinant-index identity over \(\mathbb Z\) proves the
displayed count.

The normalized pullback chain rule contributes
\(\operatorname{Nm}_a(k)=k^{d_a}\).
The normalizing polynomial contributes \(k^{\deg P}\).
They, and \(p^{v_p(I_N)}/I_N\), are units in \(V\).
After cancelling these units the term for this prime-to-primary
factor is a unit times
\(s\prod_\tau\tau(\det P)^{r_\tau}z\).
Its valuation is at least \(-\rho\) by
\eqref{eq:ftb-adapted-scalar-integral}.
Sum over the prime-to-primary factor and use the valuation
inequality.  No bound on the adapted \(k,k_1,\gamma\) occurs.
\end{proof}

\begin{srlemma}[The level-prime cancellation]
\label{lem:ftb-level-prime-cancellation}
Assume \(\operatorname{TrivAt}(\Lambda,p)\).
Suppose
\[
 N=\beta\alpha n',\quad (n',p)=1,\quad
 \alpha\bar\alpha=A,\quad\beta\bar\beta=B,\quad AB=p^j ,
\]
and \([\beta]t_b\) is killed by \(n'^2A\).
Then
\[
 v\!\left(\frac{\operatorname{Nm}_a(\alpha)}
                     {I_N N^{\deg P}}
            \sum_{z\in Z_\Lambda(N)}
                  \operatorname{hv}(F_{z,t},P)\right)\ge-\rho .
\]
\end{srlemma}
\begin{proof}
An integer \(k\) prime to \(p\) makes both lattice inclusions
\(k\Lambda\subset R^n\) and \(kR^n\subset\Lambda\) hold.
Their quotient indices have no \(p\)-part.
Thus \(v_p(J_N)=v_p(I_N)\).
It also follows that the determinant ideals of both \(\Lambda\)
and \(\Lambda_0\) are units above \(p\); hence
\(v(\beta_\Lambda)=0\).
The polynomial \(P\) has integral coefficients by the same
linear-generator argument with \(\gamma=1\).

For a primary \(z\in Z_\Lambda(N)\), the defining lattice maps
and \(kR^n\subset\Lambda\) imply \(Nz=0\), after cancelling \(k\).
Consequently the first normalized translate of the
multiplication-by-\(N\) pullback is unchanged by this primary
factor.  Symmetric translations of coprime orders commute.
In the product decomposition \(Z_p\times Z^{(p)}\), the harmonic
value therefore depends only on the second factor.
The full sum is exactly \(\#Z_p\) times that smaller sum.
As \(v_p(\#Z_p)=v_p(I_N)\), the quotient \(\#Z_p/I_N\) is a unit.

The prime-to-\(p\) factor \(z'\), being killed by \(N^2=n'^2p^j\),
is killed by \(n'^2\).
The mixed translate-pullback identity of
Proposition~\ref{prop:ftb-mixed-value-bound} now applies with these
\(\alpha,\beta,n'\).
It bounds each remaining term after multiplication by
\(\operatorname{Nm}_a(\alpha)N^{-\deg P}\).
Summing, and cancelling the unit \(\#Z_p/I_N\), proves the result.
\end{proof}

In the same-model application the lift of a valued framed point
has an exact theta-value formula
\begin{equation}\label{eq:ftb-local-theta-value}
 \beta_\Lambda g
  =\frac{\zeta}{C_M}
        \sum_q w_q
         \left(\frac1{I_N N^{\deg P}}
               \sum_z\operatorname{hv}(F_{z,t(q)},P)\right),
 \qquad
 C_M=\#(\operatorname{LevQ}^n)=M^{frn}.
\end{equation}
Here \(w_q\) is an algebraic integer and \(\zeta\) is a root of
unity; for the identity cover it is exactly \(1\).
The division point satisfies
\(t_b(q)^N=\operatorname{levPt}(q_b)\) and is killed by some power
of the half-level ideal \(\mathfrak M\).
This equality is on the same chart after a finite extension of
complete DVRs.  If its ramification factor is \(e'\), its
valuation is \(e'v\); every normalized inequality divides by the
same \(e'\), so the extension causes no loss.
The existence of this lift with its \(\mathrm{Good3}\) chart is
the precise geometric hypothesis
\srcid{EvROfFrMapFromTest.short\_value\_on\_cover}.
Its geometric producers are
\srcid{ExistsFiniteExtensionIntegralFramedPointLiftCove} and
\srcid{ExistsFiniteExtUniformisedChartRationalTorsion}; these must
be matched to the root lane's actual N3/model construction.
The formula itself follows by naturality of framed evaluations,
faithfully flat theta-value descent, and the same coefficient-field
embedding; it is not an assertion about a different model lattice.

Let \(d=d_a\), \(\Delta=|\operatorname{disc}E|\), and
\begin{equation}\label{eq:ftb-one-Z}
                       Z=C_M\Delta^dM^d,\qquad N_a=Z^2 .
\end{equation}
The following is the exact local closure used for the integral
theta family.

\begin{srproposition}[One denominator at every finite place]
\label{prop:ftb-one-denominator-all-places}
Under the same-model chart, reached-section and formal comparison
conditions above, let \(y\) be a raw theta section from
\(\Lambda,N,P,x_0\), and let
\(y_1=sy\) satisfy \eqref{eq:ftb-small-scalar-ideal}.
Assume \(\operatorname{TrivAt}(\Lambda,M)\).
Then, for every finite-place valuation in the integral-section
lattice definition,
\[
                              v(Z^2y_1)\ge0 .
\]
The integer \(Z\) is independent of \(\Lambda,N,P,x_0,s\) once the
row degrees, rank and level are fixed.
In the notation of the size hypothesis,
\[
                         \log N_a\le134(frL)^2\kappa .
\]
\end{srproposition}
\begin{proof}
Put \(g=\beta_\Lambda^{-1}y\); this is the exact identity-cover
short rule corresponding to the raw theta section.
Thus \eqref{eq:ftb-local-theta-value} has left side \(y\).

If \(p\nmid M\), the half-level-primary division point \(t(q)\)
has order prime to \(p\).
Indeed it is killed by a power of \(\mathfrak M\), hence by a
power of the integer \(M\).
Lemma~\ref{lem:ftb-coprime-translates-scalar}, applied separately
to each \(q\), shows
\[
                         v(sy)\ge-v(C_M)-\rho .
\]
For the host,
\(\rho=(d/2)\sum_{l\mid\operatorname{cond}E}
 v(l)\operatorname{ord}_l\Delta/[E:\mathbb Q]\).
Only \(l=p\) can contribute.  In particular
\(\rho\le d\,v(\Delta)\).
The factors \(C_M\Delta^d\) in \(Z\) clear both losses;
the factors \(M^d\) and the extra \(Z\) have nonnegative
valuation.

Suppose \(p\mid M\) but \(p^2\nmid M\).
Write \(N=p^e n'\), \((n',p)=1\).
Set \(\alpha=p\), \(\beta=p^{e-1}\), so
\(A=p^2\), \(B=p^{2e-2}\), \(AB=p^{2e}\).
The level-point equation says that \(t(q)^N\) is killed by \(M\),
so \(t(q)\) is killed by \(NM\).
Write \(M=pm\); from \(M\mid N\) we have \(m\mid n'\).
Since
\(\beta n'^2p^2=p^{e+1}n'^2\) is a multiple of
\(NM=p^{e+1}n'm\), the point \([\beta]t(q)\) is killed by
\(n'^2p^2\).
Apply Lemma~\ref{lem:ftb-level-prime-cancellation}.
Since \(\operatorname{Nm}_a(p)=p^d\), it gives
\[
                         v(y)\ge-v(C_M)-d\,v(p)-\rho .
\]
As \(p\mid M\), the \(M^d\) factor clears this additional loss.

For completeness, if \(p^2\mid M\), write \(M=p^km\),
\(N=p^e n'\), \(p\nmid mn'\), \(k\le e\).
Use \(\alpha=p^k\), \(\beta=p^{e-k}\);
then \(A=p^{2k}\), \(B=p^{2(e-k)}\).
The same torsion equations show
\([\beta]t(q)\) is killed by \(n'^2p^{2k}\).
More explicitly, \(t^N\) is killed by \(M\), and
\(m\mid n'\) follows from \(M\mid N\).
The exponent
\(p^{e-k}n'^2p^{2k}=p^{e+k}n'^2\)
is therefore a multiple of
\(NM=p^{e+k}n'm\).
This proves the assertion directly on the actual torsion point.
The lemma gives the same bound with \(kd\,v(p)\);
as \(p^k\mid M\), \(M^d\) clears it.
This is the source's separate \srcid{factor\_pow} proof; it does
not require \(M\) squarefree.

At primes dividing \(M\), the scalar \(s\) itself has nonnegative
valuation.  In fact \(\operatorname{TrivAt}(\Lambda,p)\)
makes both local determinant ideals units, so
\eqref{eq:ftb-small-scalar-ideal} forces \(v(s)\ge0\).
The preceding two cases thus bound \(v(sy)\) as well.
The nonarchimedean inequality, the integrality of every \(w_q\),
and \(v(\zeta)=0\) justify summing \eqref{eq:ftb-local-theta-value}.
They prove \(v(Zy_1)\ge0\) in each local case where this intermediate
bound is printed; the source records the uniform conclusion with
\(Z^2\).
Multiplying any section of nonnegative valuation by the positive
integer \(Z\) preserves that inequality, either by the scalar
valuation formula or by repeated addition.
Hence the displayed source conclusion follows.

The lattice of integral cusp sections is the intersection of
these valuation conditions inside the same characteristic-zero
cusp space.
Thus \(Z^2y_1\) belongs to it.
The calculation of
\(\log Z\le67(frL)^2\kappa\) was proved in the denominator-size
lemma; taking the square gives \(134\).
No choice made at an individual prime alters \(Z\).
\end{proof}

\paragraph{Source concordance and geometric boundary.}
All paths in this paragraph are under
\srcid{src/Development/AbelianVarietyTheory/Proofs/} at the frozen
source tag.
The adapted construction is
\srcid{RingsOfIntegers/ExistsAdaptedSimilitudeDataAtPrime.lean}
(\srcid{1b6eb2a2be303af09b0fd1507f70b85103ead2a0}),
using \srcid{ExistsUnitaryOfSelfDual.lean}
(\srcid{6eb2d1c33613ce009b53f1112abd0b4f5812474d})
and \srcid{IsPrincipalIdealRingLoc.lean}
(\srcid{1ad5550fe752c799883e6ebec4eac2e06291c2e7}).
Its Cayley proof is
\srcid{UnitaryExistence/FiniteNotIsUnitOneAddSmul.lean}
(\srcid{e7018a2755c0283f6ca79dac084cd681978e09e5}).
The determinant and index proofs are
\srcid{RingsOfIntegers/IndexRangeMatrix.lean}
(\srcid{0d0ef8c8599afc5bd8f7f818e7428398b008b4ad}),
\srcid{CardPrimaryTorsionTuplesKilledByMatrix.lean}
(\srcid{067d9220f5d5567272a1ab2b5adf1532beb6c248}),
and
\srcid{DivisorProofs/IntersectedLatticeDetIdealTypeNormIntegral.lean}
(\srcid{05cc2d038e4d187a740c9036f2de72f585fe1890}).
The exact kernel criterion is
\srcid{RingsOfIntegers/InZhIffKillsAdapted.lean}
(\srcid{2b2f3604c14f57d3a8f6821ba0c6f803a55594e0}).
The coefficient substitution is
\srcid{RingsOfIntegers/MemRangeOfIsIntegral.lean}
(\srcid{dd3360a4d7aa612b89e83c1cd91707f35cc7f6cf}).

The finite-coset proof is \srcid{RingsOfIntegers/SumHvPrimaryKernel.lean};
its model inputs are kept exactly as in the preceding formal-kernel
translation.  The scalar refinement is
\srcid{DivisorProofs/ShortSeriesBoundCoprimeTorsionMul.lean}
(\srcid{17adcd0e29a6cea26cde4903bc003b21eb3002cd}).
The trivial-lattice determinant/index proof is
\srcid{RingsOfIntegers/LatticeTrivialAtPrimeIndicesAndUnits.lean}
(\srcid{2e5f73cdd3ffcbced2cb3c5a5b4e461c4f002579}).
The three local applications are
\srcid{DivisorProofs/ClearingIntegerShortSeriesIntegralCoprimeLevel.lean},
\srcid{SmulIntZNonnegAway.lean}, and
\srcid{SmulIntZSqNonnegOfSqDvd.lean}
(\srcid{ee0c23239d1d51bb4aaa76ee0f7e8c69d47812b9}).
Their final assembly is
\srcid{DivisorProofs/ExistsBoundedIntegerGoodCoverThetaSeriesIntegral.lean},
\srcid{integral\_smul\_thLa}.
The finite-extension comparison is
\srcid{RingsOfIntegers/EvROfFrMapFromTest.lean}
(\srcid{26d9b6a33485f0f314af3d908366bc56aa158067}).
This subsection proves its numerical, determinant, torsion-count
and finite-coset consequences.  It does not replace the actual
same-model lift, integral theta-group action, formal gauge,
uniformization or expansion constructors by generic extension
principles.
